% GENERATED FILE. Edit canonical lessons, then run npm run generate:books.
% canonical-source-hash: fnv1a64:a4a9bbb475cdce93
% canonical-lessons: GU-C375-avasyak, GU-C375-anivarya, GU-C375-pharjiyat, GU-C375-patvu, GU-C375-puru-padvu

\chapter{What Must Come First}
\label{ch:gu-what-must-come-first}
\hlchaptermodality{375}

\begin{chapteropening}
\textbf{By the end of this chapter:} \emph{I can say what is needed before something can happen.}

Complete the last lesson of chapter 375: I can say what is needed before something can happen.
\end{chapteropening}

\section[\texorpdfstring{\gu{આવશ્યક} (āvaśyak)}{આવશ્યક (āvaśyak)}]{\gu{આવશ્યક} (āvaśyak) — essential}

\label{lesson:GU-C375-avasyak}



\begin{quote}
Before the new one: say the Gujarati for to keep in check, then the Gujarati for to avert.
\end{quote}

\subsection*{You'll want to know: \gu{આવશ્યક}}
\textbf{\gu{આવશ્યક}} — \emph{āvaśyak} — \textquotedblleft{}essential\textquotedblright{}.

\begin{quote}
\gu{જો} \gu{તમે} \gu{બેંકમાં} \gu{ખાતું} \gu{ખોલશો} \gu{તો} \gu{ઓળખપત્ર} \gu{આવશ્યક} \gu{છે}.
\end{quote}

\begin{quote}
\emph{jo tame beṅkmā̃ khātũ kholśo to oḷakhpatra āvaśyak che.}
\end{quote}

\begin{quote}
If you open a bank account, an identity card is essential.
\end{quote}

\begin{cousinweb}
\textbf{\gu{આવશ્યક}} (\emph{āvaśyak}) is a tatsama, borrowed from Sanskrit \emph{āvaśyaka}, \textquotedblleft{}necessary, inevitable\textquotedblright{} (Monier-Williams).
\end{cousinweb}

\begin{grammarlens}[title={What you've built}]
One more word for saying what happens if.
\end{grammarlens}

\subsection*{Your turn}
\begin{itemize}
\raggedright
  \item \emph{Say it:} \emph{āvaśyak}
  \item \emph{Say it:} \emph{āvaśyak}, once more
  \item \emph{Say it:} \emph{jo tame beṅkmā̃ khātũ kholśo to oḷakhpatra āvaśyak che.} — If you open a bank account, an identity card is essential.
  \item \emph{Recall:} say the Gujarati for to keep in check, then the Gujarati for to avert, then say \emph{āvaśyak} again
\end{itemize}

\begin{tcolorbox}[breakable,colback=teal!4,colframe=teal!35!black,title={Before you move on}]
What does \gu{આવશ્યક} mean? (\textquotedblleft{}Essential\textquotedblright{}.) Say it once more.
\end{tcolorbox}
\section[\texorpdfstring{\gu{અનિવાર્ય} (anivārya)}{અનિવાર્ય (anivārya)}]{\gu{અનિવાર્ય} (anivārya) — unavoidable}

\label{lesson:GU-C375-anivarya}



\begin{quote}
Before the new one: say the Gujarati for to avert, then the Gujarati for essential.
\end{quote}

\subsection*{You'll want to know: \gu{અનિવાર્ય}}
\textbf{\gu{અનિવાર્ય}} — \emph{anivārya} — \textquotedblleft{}unavoidable\textquotedblright{}.

\begin{quote}
\gu{જો} \gu{પેટ્રોલનો} \gu{ભાવ} \gu{વધશે} \gu{તો} \gu{બસનું} \gu{ભાડું} \gu{વધવું} \gu{અનિવાર્ય} \gu{છે}.
\end{quote}

\begin{quote}
\emph{jo peṭrolno bhāv vadhśe to basnũ bhāḍũ vadhvũ anivārya che.}
\end{quote}

\begin{quote}
If the price of petrol goes up, a rise in bus fares is unavoidable.
\end{quote}

\begin{cousinweb}
\textbf{\gu{અનિવાર્ય}} (\emph{anivārya}) is a tatsama, borrowed from Sanskrit \emph{anivārya}, \textquotedblleft{}not to be warded off, unavoidable\textquotedblright{} (Monier-Williams). It is built on the same root as \textbf{\gu{નિવારવું}} (\emph{nivārvũ}), \textquotedblleft{}to avert\textquotedblright{}, with \emph{a-}, \textquotedblleft{}not\textquotedblright{}.
\end{cousinweb}

\begin{grammarlens}[title={What you've built}]
One more word for saying what happens if.
\end{grammarlens}

\subsection*{Your turn}
\begin{itemize}
\raggedright
  \item \emph{Say it:} \emph{anivārya}
  \item \emph{Say it:} \emph{anivārya}, once more
  \item \emph{Say it:} \emph{jo peṭrolno bhāv vadhśe to basnũ bhāḍũ vadhvũ anivārya che.} — If the price of petrol goes up, a rise in bus fares is unavoidable.
  \item \emph{Recall:} say the Gujarati for to avert, then the Gujarati for essential, then say \emph{anivārya} again
\end{itemize}

\begin{tcolorbox}[breakable,colback=teal!4,colframe=teal!35!black,title={Before you move on}]
What does \gu{અનિવાર્ય} mean? (\textquotedblleft{}Unavoidable\textquotedblright{}.) Say it once more.
\end{tcolorbox}
\section[\texorpdfstring{\gu{ફરજિયાત} (pharjiyāt)}{ફરજિયાત (pharjiyāt)}]{\gu{ફરજિયાત} (pharjiyāt) — compulsory}

\label{lesson:GU-C375-pharjiyat}



\begin{quote}
Before the new one: say the Gujarati for essential, then the Gujarati for unavoidable.
\end{quote}

\subsection*{You'll want to know: \gu{ફરજિયાત}}
\textbf{\gu{ફરજિયાત}} — \emph{pharjiyāt} — \textquotedblleft{}compulsory\textquotedblright{}.

\begin{quote}
\gu{જો} \gu{તમે} \gu{આ} \gu{શાળામાં} \gu{ભણશો} \gu{તો} \gu{ગણવેશ} \gu{ફરજિયાત} \gu{છે}.
\end{quote}

\begin{quote}
\emph{jo tame ā śāḷāmā̃ bhaṇśo to gaṇveś pharjiyāt che.}
\end{quote}

\begin{quote}
If you study at this school, the uniform is compulsory.
\end{quote}

\begin{cousinweb}
\textbf{\gu{ફરજિયાત}} (\emph{pharjiyāt}) is built on \textbf{\gu{ફરજ}} (\emph{pharaj}), \textquotedblleft{}a duty\textquotedblright{}, usually traced to Arabic \emph{farḍ}, \textquotedblleft{}obligation\textquotedblright{}, with the ending \emph{-iyāt}: what duty requires.
\end{cousinweb}

\begin{grammarlens}[title={What you've built}]
One more word for saying what happens if.
\end{grammarlens}

\subsection*{Your turn}
\begin{itemize}
\raggedright
  \item \emph{Say it:} \emph{pharjiyāt}
  \item \emph{Say it:} \emph{pharjiyāt}, once more
  \item \emph{Say it:} \emph{jo tame ā śāḷāmā̃ bhaṇśo to gaṇveś pharjiyāt che.} — If you study at this school, the uniform is compulsory.
  \item \emph{Recall:} say the Gujarati for essential, then the Gujarati for unavoidable, then say \emph{pharjiyāt} again
\end{itemize}

\begin{tcolorbox}[breakable,colback=teal!4,colframe=teal!35!black,title={Before you move on}]
What does \gu{ફરજિયાત} mean? (\textquotedblleft{}Compulsory\textquotedblright{}.) Say it once more.
\end{tcolorbox}
\section[\texorpdfstring{\gu{પતવું} (patvũ)}{પતવું (patvũ)}]{\gu{પતવું} (patvũ) — to get done}

\label{lesson:GU-C375-patvu}



\begin{quote}
Before the new one: say the Gujarati for unavoidable, then the Gujarati for compulsory.
\end{quote}

\subsection*{You'll want to know: \gu{પતવું}}
\textbf{\gu{પતવું}} — \emph{patvũ} — \textquotedblleft{}to get done\textquotedblright{}.

\begin{quote}
\gu{જો} \gu{કામ} \gu{વહેલું} \gu{પતશે} \gu{તો} \gu{આપણે} \gu{ફિલ્મ} \gu{જોવા} \gu{જઈશું}.
\end{quote}

\begin{quote}
\emph{jo kām vahelũ patśe to āpaṇe philm jovā jaīśũ.}
\end{quote}

\begin{quote}
If the work gets done early, we will go and see a film.
\end{quote}

Past (it …): \gu{તે} \gu{કામ} \gu{પત્યું}, \emph{te kām patyũ}.

\begin{cousinweb}
\textbf{\gu{પતવું}} (\emph{patvũ}) is the intransitive partner of \textbf{\gu{પતાવવું}} (\emph{patāvvũ}), \textquotedblleft{}to finish\textquotedblright{}; CDIAL lists the Gujarati verb, \textquotedblleft{}to be settled\textquotedblright{}, under Sanskrit \emph{prāpta}, \textquotedblleft{}obtained\textquotedblright{} (CDIAL 8945).
\end{cousinweb}

\begin{grammarlens}[title={What you've built}]
One more word for saying what happens if.
\end{grammarlens}

\subsection*{Your turn}
\begin{itemize}
\raggedright
  \item \emph{Say it:} \emph{patvũ}
  \item \emph{Say it:} \emph{patvũ}, once more
  \item \emph{Say it:} \emph{jo kām vahelũ patśe to āpaṇe philm jovā jaīśũ.} — If the work gets done early, we will go and see a film.
  \item \emph{Recall:} say the Gujarati for unavoidable, then the Gujarati for compulsory, then say \emph{patvũ} again
\end{itemize}

\begin{tcolorbox}[breakable,colback=teal!4,colframe=teal!35!black,title={Before you move on}]
What does \gu{પતવું} mean? (\textquotedblleft{}To get done\textquotedblright{}.) Say it once more.
\end{tcolorbox}
\section[\texorpdfstring{\gu{પૂરું} \gu{પાડવું} (pūrũ pāḍvũ)}{પૂરું પાડવું (pūrũ pāḍvũ)}]{\gu{પૂરું} \gu{પાડવું} (pūrũ pāḍvũ) — to provide}

\label{lesson:GU-C375-puru-padvu}



\begin{quote}
Before the new one: say the Gujarati for compulsory, then the Gujarati for to get done.
\end{quote}

\subsection*{You'll want to know: \gu{પૂરું} \gu{પાડવું}}
\textbf{\gu{પૂરું} \gu{પાડવું}} — \emph{pūrũ pāḍvũ} — \textquotedblleft{}to provide\textquotedblright{}.

\begin{quote}
\gu{જો} \gu{સરકાર} \gu{પાણી} \gu{પૂરું} \gu{પાડશે} \gu{તો} \gu{ખેતરો} \gu{લીલાં} \gu{થશે}.
\end{quote}

\begin{quote}
\emph{jo sarkār pāṇī pūrũ pāḍśe to khetro līlā̃ thaśe.}
\end{quote}

\begin{quote}
If the government supplies water, the fields will turn green.
\end{quote}

Past, with \gu{તેણે} (he or she …): \gu{તેણે} \gu{બધું} \gu{સાધન} \gu{પૂરું} \gu{પાડ્યું}, \emph{teṇe badhũ sādhan pūrũ pāḍyũ}.

\begin{cousinweb}
\textbf{\gu{પૂરું} \gu{પાડવું}} (\emph{pūrũ pāḍvũ}) is built from \emph{pūrũ}, \textquotedblleft{}full\textquotedblright{}, which CDIAL traces to Sanskrit \emph{pūra}, \textquotedblleft{}filling\textquotedblright{} (CDIAL 8330), and \emph{pāḍvũ}, \textquotedblleft{}to make fall\textquotedblright{}, the causative of \textbf{\gu{પડવું}} (\emph{paḍvũ}), from Sanskrit \emph{patati} (CDIAL 7722): to make a supply full.
\end{cousinweb}

\begin{grammarlens}[title={What you've built}]
One more word for saying what happens if.
\end{grammarlens}

\subsection*{Your turn}
\begin{itemize}
\raggedright
  \item \emph{Say it:} \emph{pūrũ pāḍvũ}
  \item \emph{Say it:} \emph{pūrũ pāḍvũ}, once more
  \item \emph{Say it:} \emph{jo sarkār pāṇī pūrũ pāḍśe to khetro līlā̃ thaśe.} — If the government supplies water, the fields will turn green.
  \item \emph{Recall:} say the Gujarati for compulsory, then the Gujarati for to get done, then say \emph{pūrũ pāḍvũ} again
\end{itemize}

\begin{tcolorbox}[breakable,colback=teal!4,colframe=teal!35!black,title={Before you move on}]
What does \gu{પૂરું} \gu{પાડવું} mean? (\textquotedblleft{}To provide\textquotedblright{}.) Say it once more.
\end{tcolorbox}
